\documentclass[10pt,a4paper]{article}
\usepackage[margin=19mm]{geometry}
\usepackage{amsmath,amssymb,amsthm,booktabs,microtype}
\usepackage[colorlinks=true,urlcolor=blue,citecolor=blue,linkcolor=blue]{hyperref}
\usepackage{enumitem}
\setlist{nosep,leftmargin=*}
\setlength{\parskip}{3pt}
\setlength{\parindent}{0pt}
\setlength{\emergencystretch}{2em}
\newcommand{\tht}{\vartheta}
\newcommand{\RR}{\mathbb R}
\newcommand{\EE}{\mathcal E}
\newcommand{\Raff}{R_{\rm aff}}
\newcommand{\Rold}{R_{\kappa}}
\newcommand{\Reg}{\mathcal R}
\newtheorem{lemma}{Lemma}
\title{Affine detector counting with asymmetric source savings}
\author{Zhongpai Gao\\\small With AI assistance; predecessor contributions acknowledged below}
\date{15 October 2026}
\begin{document}
\maketitle
\vspace{-2.5em}
\begin{center}\small\textbf{Official tracker verification accepted; PR60 merged 15 October 2026.}\end{center}
\begin{abstract}
We prove nonvanishing of every Dirichlet
$L$-function in $\Re s>21871/25000=0.874840$, away from the principal pole.
The construction combines timbalud's asymmetric marked HIGH count and
common-radical saving with DaysSky's use of Collet's initial affine moment
inside the detector count. The stronger source moment increases the detector's
available prime-slot width. A small change of the height junction preserves
the actual source gates. We give the counting formulas, the allocation
argument, and the passage from uniform scalar inequalities to the inherited
analytic continuation argument. Exact rational certificates cover the continuous parameter ranges, and the
complete source assembly has passed canonical comparison and all three kernels.
\end{abstract}

\section{Framework and theorem}
Let $\beta\leq7/8$ be the supremum of the nontrivial zero real parts in
Cycle61's closed Eisenstein Hecke family. The upper bound and the transfer from
this family to all Dirichlet characters are inherited from the OpenAI and
Collet framework \cite{collet,openai}. Set
\begin{gather}
 \tht=\frac{21871}{25000},\quad h=\frac{1989}{2500},\quad
 d_-=\frac{159}{200},\quad d_+=h+10^{-10},\quad
 \kappa_*=2\tht-1=\frac{9371}{12500},\label{parameters}\\
 \ell=\frac{11}{3}-4\tht=\frac{3137}{18750},\quad
 x=1+\ell-h=\frac{13939}{37500},\quad
 y=h-2\ell=\frac{17287}{37500},\quad
 \lambda=\frac{19312}{4787},\quad
 B=\frac{14200}{14361}+\frac7{10000}.\nonumber
\end{gather}
\textbf{Theorem.} For every modulus $q\geq1$, every Dirichlet character
$\chi$ modulo $q$, and every $s\in\mathbb C$,
\[
 \tht<\Re s,\qquad \neg(\chi=1\ \text{and}\ s=1)
 \quad\Longrightarrow\quad L(\chi,s)\ne0.
\]
The accompanying Lean construction proves this canonical statement.

Assume $\beta>\tht$, keeping $\beta$ as the global supremum. The continuation
contract first chooses a common error reserve $\sigma>0$ and
$\omega=(\beta-\tht)/2$, before quantifying over primitive targets.
Constants and a small height exponent may depend on the target character;
the resulting positive continuation gap must remain uniform. This order
is essential when the affine moment is applied to detector characters.

\section{A stronger actual source moment}
Collet's initial affine moment bounds the normalized plain energy, including
its selected prime factors, by
\begin{equation}
 \EE_\eta(U;m,s)\ll_{\eta,\varepsilon}
 (1+T)^{D_\eta}\,\mathcal S(\Phi)^2 U^{1+\varepsilon},
 \qquad 2m+\lambda s\leq B.\label{affine-moment}
\end{equation}
Here $m$ is the plain length, $s$ the selected logarithmic prime width,
$T$ controls the shifts, and $\mathcal S$ is a finite profile seminorm.
All fixed masks, relative characters and source normalizations are those of
the actual initial moment. We use \eqref{affine-moment} only in its proved
window, with $s\leq1/25$ and $m$ in the interval specified in Section~\ref{sec:detector}.

The admissible choice of $B$ is justified by the complete marked source
argument of timbalud \cite{asymmetric}, not by an additional moment hypothesis.
The useful common-allocation step can be seen directly. For first-allocation
coordinates $c,d$, radical coordinate $R$, cancellation coordinate $R_0$ and
retained width $L_c$, the actual source contains
\begin{equation}
 F=(1-\rho)(c-d)_+-\delta c-(1+\rho)R-E
       -\frac{(3c-5d-R)_+}{6}+\rho L_c+R_0.\label{common-word}
\end{equation}
The radical penalty is tied to the same $c,d,R$; it cannot be discarded and
then varied independently. Under the literal gates
$R\leq\min(c,d)$, $L_c\leq5(R-R_0)$, $E,R_0,c,d\geq0$ and
$1/4\leq\rho\leq2/5$, one has
\begin{equation}
 F\leq-(\delta-3\rho/2)c-E.\label{coupled}
\end{equation}
Indeed, for $d\leq c$, write $a=2-5\rho\in[0,1]$ and use
$a d+(c/2-d)_+\geq ac/2$. For $c\leq d$, use $R\leq c$ and
$1-5\rho/2\geq0$. These two cases give \eqref{coupled} on the whole domain.

In the marked HIGH range $17/25\leq\delta\leq7/10$,
$\rho=1-10\delta/[3(5-2\delta)]\leq103/273$; thus the first removed
width has rate at least $a_0=519/4550$. The second allocation has the
stronger direct rate $\delta$. With $M$ the source scale, let $u,f,v\geq0$
be the common, residual frozen and surviving widths, and let $Q$ be the
actual remaining frequency scale. The source identities and two-stage
collision accounting give
\[
 u+f+v\geq\frac{33M}{1000},\qquad
 Q\geq\frac{1121M}{1250}-2u.
\]
Asymmetric inverse/plain selections supply the live-width payment, so that
\begin{equation}
 a_0u+\delta f+\frac{224}{1755}\min\!\left(v,\frac{13Q}{400}\right)
 \ \geq\ \frac{37M}{10000}.\label{ledger}
\end{equation}
For $v\leq13Q/400$, use the width supply and the smallest of the three
rates. In the other case insert the lower bound for $Q$; the coefficient
of $u$ remains positive and the constant term already exceeds the right
side. This proves \eqref{ledger} without choosing an allocation extremizer.
The complementary large sixth-power tail is paid by its inherited bound.
Charging at most $M/20000$ leaves $73M/20000$. These savings are retained
inside the actual source moment, with moving common ideals, frozen
coefficients and encoded-frequency collisions unchanged.

\section{Applying that moment inside the detector}\label{sec:detector}
A detector fiber has inverse length $r$, plain length $m$, cutoff $t$,
zero density $\delta$ and normalized prime amplitude $p\in[0,1/2]$.
Up to the explicitly charged mesh errors, $r+m=t$. The three relevant count
exponents are
\begin{align}
 I(\delta,p,r)&=1-\delta\bigl(p+(1-p)r\bigr),\nonumber\\
 A(\delta,p,m)&=1-2\delta m-2\delta p\frac{B-2m}{\lambda},\label{branches}\\
 L(\delta,t)&=1-\delta+(5/6-\delta)(t-1).\nonumber
\end{align}
The affine plain branch uses \eqref{affine-moment} at width
$(B-2m)/\lambda$, with a small subtraction for rounding. It improves on
the old plain width $(1-2m)/(6\kappa)$ on the critical fibers.

This is a legitimate instance of the \emph{actual} source theorem.
Fix a final target character $\psi$, and write $\eta_j$ for its detector
dictionary characters. Each $\eta_j$ is represented as its own natural row,
with the exact excluded-prime support and fixed masks. The generic
paid main and error families apply to that row, and the literal energy
identity identifies their energy with the fiber energy. The marked HIGH
mesh cap is intersected with the original analytic mesh before the slot
partition is chosen. This dictionary is finite: we take the maximum of its
moment degrees and the union of its seminorm sets, then choose a common
small height exponent for all its labels. The resulting exponent
$\tau_\psi$ may depend on the final target $\psi$. No affine moment or
uniform degree over infinitely many target characters is assumed.

Put
\begin{align*}
 D(p)&=3-p-4p/\lambda,& a(p)&=2-4p/\lambda,& c(p)&=1-2pB/\lambda,\\
 J(\delta,p)&=(5/6-\delta)D(p)+\delta(1-p)a(p).
\end{align*}
Solving $I=A$ at $m=t-r$ and then balancing with $L$ gives
\begin{align}
 r_N(p,t)&=\frac{a(p)t+2pB/\lambda-p}{D(p)},\qquad
 t_N=1+\frac{\delta(1-p)c(p)}{J(\delta,p)},\nonumber\\
 \Raff(\delta,p)&=1-\delta+
 \frac{(5/6-\delta)\delta(1-p)c(p)}{J(\delta,p)}.\label{new-count}
\end{align}
Use this count on
\[
 \Reg=[89/200,97/200]\times[199/400,1/2],\qquad
 m\in[m_0,847/2000],\quad m_0=(B-\lambda/25)/2.
\]
The exact region certificates give $D,J>0$, $1\leq t_N\leq3/2$,
$23/37\leq r_N(p,t_N)\leq1$, and
$t_N-r_N(p,t_N)\geq m_0+10^{-5}$.
For selected widths below $17/500$, the old $\kappa$-marked moment is used;
the affine theorem is invoked only at or above this lower gate. For plain lengths
above $847/2000$, the old marked plain branch gives the same short exponent.
The original unmarked branches cover the remaining cases. Thus the complete
fiber argument gives \eqref{new-count}, with the original charged losses.
The detector spacing is chosen $\varepsilon\leq10^{-5}$; the actual baseline
constructor already gives $\varepsilon\leq1/200000$.

\section{Closing every row and returning to the $L$-function}
For $\kappa\in[\kappa_*,3/4]$, write
\[
 C_\kappa(p)=\frac{2(1-p)(3\kappa-p)}{3\kappa(3-p)-2p},\qquad
 \Rold(\delta,p)=1-\delta+
 \frac{\delta C_\kappa(p)(5/6-\delta)}{2(5/6-\delta+\delta C_\kappa(p))}.
\]
Let $R^*_{\kappa}$ equal $\Raff$ on $\Reg$ and $\Rold$ elsewhere.
The affine region count is at most the old count at $\kappa_*$, and
$R^*_{\kappa}$ is nonincreasing in $p$. This monotonicity is used when
actual amplitude classes are transferred through the numerator masks.

The complete normalized reconstruction has row exponent
\begin{equation}
 E_\kappa(d,C;\delta,p)=\frac{1+\delta}{2}+\frac{13h}{75}
 -\frac y2-\frac\ell2+\ell\delta p
 +d\left(\frac14+\frac34R^*_\kappa(\delta,p)-\frac{17}{50}\right)
 -\frac{\delta p C}{2}.\label{row}
\end{equation}
Its two physical capacities are
\[
 C_{\rm new}(d)=\frac{2y-(2-B)d}{\lambda},\qquad
 C_{{\rm old},\kappa}(d)=\frac{\min(d,2y-d)}{6\kappa}.
\]
With \eqref{parameters}, rational Bernstein certificates on closed
rectangles, together with affine interpolation in $d$, establish
\begin{align}
 E_{\kappa_*}(d,C_{\rm new}(d);\delta,p)
  &<\tht-\frac{4999}{5000000000},&&d_-\leq d\leq d_+,\label{new-row}\\
 E_{\kappa_*}(d,C_{{\rm old},\kappa_*}(d);\delta,p)
  &<\tht-\frac1{100000},&&1/200\leq d\leq d_- .\label{old-row}
\end{align}
Both inequalities hold for \emph{every} $\delta\in[0,3/4]$ and
$p\in[0,1/2]$. Below $d=1/2$ the old saving is at least $1/1000$.
For clarity, the new top row uses the affine count on $\Reg$ and the old
count on its complement. The old half-height and junction rows are
certified on the whole rectangle. Positive rational Bernstein
coefficients certify the full polynomials, rather than samples of them;
the corresponding cleared denominators are positive. The extension from
$h$ to $h+10^{-10}$ costs less than $2\cdot10^{-10}$, giving
\eqref{new-row}. At $d_-$ the new capacity exceeds the old capacity,
so the new junction inherits \eqref{old-row}.

The junction $d_-=h-3/5000$ also respects the actual source gates.
In particular the plain scale $n=1-y/d+t$, for $0\leq t\leq1/1000$,
satisfies $41/100<n<264/625$; its affine capacity lies in
$(17/500,37/1000)$; and
$0\leq2y/d-1<4/25$. The fixed source HIGH, LOW, amplifier and tail
reserves, as well as the cube-floor payment, therefore remain available
from the inherited complete estimates. These gates are checked alongside
the row certificates, not inferred from their positivity.

Finally set $\kappa=2\beta-1$. The explicit feedback estimates strengthen
\eqref{new-row}--\eqref{old-row} to
\begin{align*}
 E_{2\beta-1}(d,C_{\rm new};\delta,p)-\beta
 &<-\frac{4999}{5000000000}-\frac{24}{25}(\beta-\tht),\\
 E_{2\beta-1}(d,C_{{\rm old},2\beta-1};\delta,p)-\beta
 &<-\frac1{100000}-\frac34(\beta-\tht).
\end{align*}
The finite reconstruction, mesh and height costs can consequently be
chosen below a fixed fraction of $\beta-\tht$. The unchanged source masks,
prime products, transport and regularizer yield a common error reserve
$\sigma>0$. Assembly fixes this reserve and $\omega=(\beta-\tht)/2$
before choosing each target's height. Thus the Mellin continuation gap is
\[
 \varepsilon_{\rm cont}=\min\!\left(\frac{\beta-\tht}{2},\sigma\right)>0,
\]
independent of the primitive target. Only the constants, regularizer and
height choice depend on that target. The inherited continuation theorem
excludes every primitive nontrivial zero in
$\Re s>\beta-\varepsilon_{\rm cont}$, contradicting the global supremum.
Hence $\beta\leq\tht$, and the Dirichlet transfer proves the theorem.

\section{Scope, provenance and verification}
The mathematical advance is the use of a stronger, already source-paid
affine moment \emph{inside} the detector. The proof does not add the two
predecessor improvements as numerical boundary gains. The first improvement
changes the admissible moment; the second uses that moment in a separate
counting problem. Their interaction changes the crossing formulas and the
admissible prime-amplitude region, which is why the complete row
certificates and source gates must be regenerated together.

The accompanying Lean construction exposes no new analytic premise in its
final target. Its route is
\begin{gather*}
 \text{actual per-label affine energy}\ \longrightarrow\
 \text{actual detector count}\ \longrightarrow\\
 \text{finite numerator and source probe}\ \longrightarrow\
 \text{canonical Dirichlet nonvanishing}.
\end{gather*}
The final continuation contract has the order
$\exists\omega,\sigma>0\;\forall\psi\;\exists\tau_\psi$,
with the displayed common gap. For each primitive target $\psi$, its finite
detector dictionary shares $\tau_\psi$ after taking the stated finite maxima;
character-dependent heights cannot shrink the continuation gap.
The scalar argument above explains the new mechanism. The analytic
reductions are supplied by the complete inherited source theorems.

\textbf{Verification status.}
The official \href{https://gtimg.github.io/QuasiRiemannTracker/certificates/gao-affine-asymmetric-20261015.json}{tracker certificate}
records $21871/25000$ at final source commit
\href{https://github.com/GTimG/QuasiRiemannTracker/blob/3874f3e75de18ff5806dcb9519fc260e034ef0e9/submissions/gao-affine-asymmetric-20261015/src/Candidate/Main.lean}{\texttt{3874f3e75de1}}
as verified on 15 October 2026 by Lean, NanoDa and con-ron;
\href{https://github.com/GTimG/QuasiRiemannTracker/pull/60}{PR60} was merged that day.
The accompanying sanitized historical local receipt is bound to source
commit \href{https://github.com/Gaozhongpai/QuasiRiemannTracker/blob/74643447f52a2df4d05cbb23f50b9e6f2959beaa/submissions/gao-affine-asymmetric-20261015/src/Candidate/Main.lean}{\texttt{74643447f52a}}.
In that local replay, Comparator accepted the canonical statement with only
\texttt{Classical.choice}, \texttt{Quot.sound} and \texttt{propext}, and all
three kernels accepted the proof. Its parent replay compiled all 897
candidate modules from source; its metadata-only follow-up reused an
authenticated local cache and ran fresh controls, canonical exports and
all three kernels. Exact rational checks cover the displayed continuous
rows and physical gates.

\textbf{Attribution.}
OpenAI supplies the protected $7/8$ framework and library. Hailey Collet
supplies Cycle25/Cycle61, the initial affine moment, and the complete
physical source and continuation architecture. The weighted numerator
uses Akash Levy's contribution and its inherited credits.
Timbalud supplies the marked HIGH count and the asymmetric/common-radical
source improvement \cite{asymmetric}. DaysSky supplies the actual
dictionary-character affine detector construction and its height order
\cite{detector}. Liqsweep/blueberryvertigo also recorded the count-side
affine idea and related arithmetic in PR52, whose Lean integration was
unfinished \cite{related}; that proposed endpoint is not used as a proved
input. Zhongpai Gao, with AI assistance, combines the source mechanisms,
chooses \eqref{parameters}, regenerates the exact certificates and carries
the changed source interfaces through the final assembly. Existing
OpenAI, Collet, Levy and other predecessor notices are retained. The
integration and manuscript are released under Apache-2.0, subject to the
retained predecessor and dependency licenses.

\begin{thebibliography}{9}\small
\bibitem{collet} Hailey Collet, \emph{Cycle61: all-Dirichlet boundary
$0.874861$}, complete source and manuscript,
\href{https://github.com/GTimG/QuasiRiemannTracker/pull/42}{QuasiRiemannTracker PR42} (2026).
\bibitem{asymmetric} timbalud, \emph{Asymmetric marked counting with coupled
radical allocation}, source at commit \texttt{8cb265dc3490},
\href{https://github.com/GTimG/QuasiRiemannTracker/pull/50}{QuasiRiemannTracker PR50} (2026).
\bibitem{detector} DaysSky, \emph{Cycle61 with the affine moment in the detector
count}, source at commit \texttt{67edfb7a1161},
\href{https://github.com/GTimG/QuasiRiemannTracker/pull/56}{QuasiRiemannTracker PR56} (2026).
\bibitem{related} liqsweep / blueberryvertigo, \emph{Collet's affine moment in the
detector count}, related arithmetic proposal,
\href{https://github.com/GTimG/QuasiRiemannTracker/pull/52}{QuasiRiemannTracker PR52} (2026).
\bibitem{openai} OpenAI, \emph{Quasi-Riemann hypothesis formalization}, protected
source snapshot \texttt{fd4aeeb2}, \url{https://github.com/OpenAI/math}.
See the retained source notices for Akash Levy's weighted numerator and the
full dependency attribution.
\end{thebibliography}
\end{document}
